\chapter{条件随机场}


条件随机场(Conditional Random Field, CRF)是用于标注和分析序列数据的一类判别式概率模型.
它在自然语言处理中广泛用于分词、词性标注、命名实体识别等任务.

从建模角度看, HMM建模的是联合概率$P(x, y)$, 而CRF直接建模条件概率$P(y|x)$.
因此, CRF更适合“已知输入序列, 预测输出标签序列”这一类监督学习问题.

\section{概率无向图模型}


条件随机场建立在概率无向图模型(probabilistic undirected graphical model)之上.
在无向图中, 结点表示随机变量, 边表示变量之间可能存在的依赖关系.

设无向图为
\[
G = (V, E)
\]
其中$V$是结点集合, $E$是边集合.
若随机变量集合记为$Y = (Y_v)_{v \in V}$, 则概率无向图模型用图结构刻画这些随机变量之间的条件独立关系.

在无向图模型中, 一个重要概念是团(clique), 即图中任意两点都相连的结点子集.
极大团(maximal clique)是不能再扩大的团.

\subsection{概率无向图的因子分解}


若联合分布$P(Y)$与无向图$G$相容, 则根据 Hammersley-Clifford 定理, 它可以分解为各个团势函数(potential function)的乘积:
\[
P(Y) = 1/Z \prod_C \psi_C (Y_C)
\]
其中:

\begin{itemize}
\item $C$表示图中的团.
\item $Y_C$表示团$C$上对应的随机变量.
\item $\psi_C(Y_C)$是定义在团上的非负函数, 称为势函数.
\item $Z$是规范化因子(partition function), 保证概率和为$1$.
\end{itemize}

规范化因子写作:
\[
Z = \sum_Y \prod_C \psi_C (Y_C)
\]

势函数本身不一定是概率, 但多个势函数相乘再经过归一化之后就得到一个合法的概率分布.
这类分解说明: 无向图模型通过局部相互作用来刻画整体概率结构.

\section{条件随机场的定义与形式}


设$X$和$Y$是随机变量, 其中$X$表示输入观测序列, $Y$表示输出标签序列.
若在给定$X$的条件下, 随机变量$Y$构成一个无向图模型, 并满足局部马尔可夫性质:
\[
P(Y_v | X, Y_{W \setminus v}) = P(Y_v | X, Y_{N_v})
\]
则称$P(Y|X)$是条件随机场.

这里:

\begin{itemize}
\item $W$表示全部结点集合.
\item $v$表示某个结点.
\item $N_v$表示与结点$v$相邻的结点集合.
\end{itemize}

这一定义说明: 在给定输入$X$时, 某个输出结点只与其邻居结点直接相关, 而与其他远处结点条件独立.

对于一般图结构, 条件随机场的条件分布写作:
\[
P(y|x) = 1/(Z(x)) \prod_C \Psi_C (y_C, x)
\]
其中
\[
Z(x) = \sum_y \prod_C \Psi_C (y_C, x)
\]

与无向图模型相比, 这里的规范化因子依赖于输入$x$, 因而记作$Z(x)$.

\subsection{线性链条件随机场}


在自然语言处理里最常用的是线性链条件随机场(linear-chain CRF).
它适用于输入和输出都按时间或位置排列的序列标注问题.

设输入序列为
\[
x = (x_1, x_2, \dots, x_n)
\]
输出标签序列为
\[
y = (y_1, y_2, \dots, y_n)
\]
则线性链CRF的图结构是一个链:

\begin{itemize}
\item 每个结点$y_i$表示位置$i$上的标签.
\item 边只连接相邻标签$y_{i-1}$和$y_i$.
\item 输入$x$作为全局已知条件.
\end{itemize}

因此, 线性链CRF的条件概率可写为:
\[
P(y|x) = 1/(Z(x)) \prod_{i=1}^n \Psi_i (y_{i-1}, y_i, x, i)
\]

这里通常约定$y_0$表示起始状态(start state).

\subsection{参数化形式}


为了便于学习参数, 通常将势函数写成指数形式:
\[
\Psi_i (y_{i-1}, y_i, x, i)
    = \exp(\sum_k \omega_k f_k (y_{i-1}, y_i, x, i))
\]
其中:

\begin{itemize}
\item $f_k(y_{i-1}, y_i, x, i)$是特征函数(feature function).
\item $\omega_k$是对应的权重参数.
\end{itemize}

代入后得到:
\[
P_{\omega} (y|x)
    = 1/(Z_{\omega} (x))
      \exp(\sum_{i=1}^n \sum_k \omega_k f_k (y_{i-1}, y_i, x, i))
\]

规范化因子为:
\[
Z_{\omega} (x)
    = \sum_y
      \exp(\sum_{i=1}^n \sum_k \omega_k f_k (y_{i-1}, y_i, x, i))
\]

这一形式说明, CRF本质上是一个对数线性模型(log-linear model)在序列标注问题上的推广.

\subsection{简化形式}


在线性链CRF中, 常把特征分为两类:

\begin{itemize}
\item 转移特征(transition feature): 描述相邻标签之间的关系.
\item 状态特征(state feature): 描述当前位置标签与输入之间的关系.
\end{itemize}

若记转移特征为$t_l(y_{i-1}, y_i, x, i)$, 权重为$\lambda_l$,
状态特征为$s_m(y_i, x, i)$, 权重为$\mu_m$,
则模型可改写为:
\[
P(y|x)
    = 1/(Z(x))
      \exp(
        \sum_{i=1}^n \sum_l \lambda_l t_l (y_{i-1}, y_i, x, i)
        +
        \sum_{i=1}^n \sum_m \mu_m s_m (y_i, x, i)
      )
\]

这种分法更贴合具体应用.
例如在词性标注中:

\begin{itemize}
\item 转移特征可以表示“前一个词性是名词, 当前词性是动词”.
\item 状态特征可以表示“当前词是\texttt{的}, 当前标签是助词”.
\end{itemize}

\subsection{矩阵形式}


为了进行高效计算, 可将线性链CRF写成与 HMM 类似的矩阵形式.

对每个位置$i$, 定义状态转移矩阵:
\[
M_i (x) = [M_i (y_{i-1}, y_i | x)]
\]
其中矩阵元素定义为:
\[
M_i (y_{i-1}, y_i | x)
    = \exp(\sum_k \omega_k f_k (y_{i-1}, y_i, x, i))
\]

则有:
\[
Z(x) = \sum_y \prod_{i=1}^n M_i (y_{i-1}, y_i | x)
\]

记标签集合大小为$m$, 并记
\[
\alpha_i = \begin{bmatrix}\alpha_i(y_1) \\ \alpha_i(y_2) \\ \vdots \\ \alpha_i(y_m)\end{bmatrix}
\]
\[
\beta_i = \begin{bmatrix}\beta_i(y_1) \\ \beta_i(y_2) \\ \vdots \\ \beta_i(y_m)\end{bmatrix}
\]

如果把起始状态$y_0$吸收到第一个位置的势函数中, 则可以把
\[
\alpha_1 = \begin{bmatrix}M_1(y_0, y_1 | x) \\ M_1(y_0, y_2 | x) \\ \vdots \\ M_1(y_0, y_m | x)\end{bmatrix}
\]
看成初始前向列向量.

之后的递推与 HMM 的矩阵形式非常相似:

\begin{itemize}
\item 前向递推通过$M_i^T$左乘前向列向量.
\item 后向递推通过$M_{i+1}$左乘后向列向量.
\item 规范化因子由最终前向量求和得到.
\end{itemize}

因此, 线性链CRF的概率计算、边缘概率计算和维特比搜索都可以归结为一系列矩阵运算.

\section{条件随机场的三个问题}


与HMM类似, CRF的核心也包括三个基本问题:

\begin{itemize}
\item 概率计算: 计算给定$x$和$y$下的概率、边缘概率及特征期望.
\item 推测状态: 求使$P(y|x)$最大的标签序列.
\item 模型学习: 用训练数据估计参数$\omega$.
\end{itemize}

\subsection{概率计算}


给定$P(y|x)$、输入序列$x$和输出序列$y$, 需要计算:

\begin{itemize}
\item 条件概率$P(y|x)$.
\item 节点边缘概率$P(y_i|x)$.
\item 边缘概率$P(y_{i-1}, y_i|x)$.
\item 特征在模型分布下的期望.
\end{itemize}

由于规范化因子$Z(x)$要对全部标签序列求和, 直接计算代价很高.
对于线性链CRF, 可以用前向-后向算法在线性时间内完成.

\subsubsection{前向-后向算法}


定义前向量:
\[
\alpha_i (y_i)
    = \sum_{y_1, \dots, y_{i-1}}
      \prod_{j=1}^i M_j (y_{j-1}, y_j | x)
\]

将它写成列向量形式:
\[
\alpha_i = \begin{bmatrix}\alpha_i(y_1) \\ \alpha_i(y_2) \\ \vdots \\ \alpha_i(y_m)\end{bmatrix}
\]

则前向算法可写成三步:

\begin{enumerate}
\item[1.] 初始化
\end{enumerate}
\[
\alpha_1 = \begin{bmatrix}M_1(y_0, y_1 | x) \\ M_1(y_0, y_2 | x) \\ \vdots \\ M_1(y_0, y_m | x)\end{bmatrix}
\]

\begin{enumerate}
\item[2.] 递推
\end{enumerate}
\[
\alpha_i = M_i^T \alpha_{i-1}, \quad i = 2,3,\dots,n
\]

定义后向量:
\[
\beta_i (y_i)
    = \sum_{y_{i+1}, \dots, y_n}
      \prod_{j=i+1}^n M_j (y_{j-1}, y_j | x)
\]

将它写成列向量形式:
\[
\beta_i = \begin{bmatrix}\beta_i(y_1) \\ \beta_i(y_2) \\ \vdots \\ \beta_i(y_m)\end{bmatrix}
\]

则后向算法可写成三步:

\begin{enumerate}
\item[1.] 初始化
\end{enumerate}
\[
\beta_n = \begin{bmatrix}1 \\ 1 \\ \vdots \\ 1\end{bmatrix}
\]

\begin{enumerate}
\item[2.] 递推
\end{enumerate}
\[
\beta_i = M_{i+1} \beta_{i+1}, \quad i = n-1, n-2, \dots, 1
\]

\begin{enumerate}
\item[3.] 终止
\end{enumerate}
\[
Z(x) = \sum_j \alpha_n (y_j) = \sum_j \alpha_1(y_j) \beta_1(y_j)
\]


\subsubsection{概率计算和期望值计算}


有了前向量和后向量之后, 就可以计算各种边缘概率.

\begin{enumerate}
\item[1.] 单点边缘概率
\end{enumerate}
\[
P(y_i | x) = (\alpha_i (y_i) \beta_i (y_i)) / Z(x)
\]

把它写成列向量形式, 可记为:
\[
p_i = 1/(Z(x)) \operatorname{diag}[\alpha_i] \beta_i
\]
其中$p_i$的第$j$个分量就是$P(y_i = y_j | x)$.

\begin{enumerate}
\item[2.] 相邻状态边缘概率
\end{enumerate}
\[
P(y_{i-1}, y_i | x)
    = (\alpha_{i-1}(y_{i-1}) M_i (y_{i-1}, y_i | x) \beta_i (y_i)) / Z(x)
\]

把所有相邻状态边缘概率写成矩阵, 记为:
\[
P_i = 1/(Z(x)) \operatorname{diag}[\alpha_{i-1}] M_i \operatorname{diag}[\beta_i]
\]
则矩阵$P_i$的第$(r, s)$个元素就是
\[
P(y_{i-1}=y_r, y_i=y_s | x)
\]

若特征函数是$f_k (y_{i-1}, y_i, x, i)$, 则其在模型分布下的期望为:
\[
E_{P_{\omega} (Y|x)}[f_k]
    = \sum_{i=1}^n \sum_{y_{i-1}, y_i}
      f_k (y_{i-1}, y_i, x, i) P(y_{i-1}, y_i | x)
\]

因此, 特征期望的计算本质上依赖于各个位置上的边缘概率矩阵$P_i$.

这个期望在参数学习中非常重要, 因为梯度恰好由“经验特征期望”减去“模型特征期望”构成.

\subsection{推测状态}


给定$P(y|x)$和输入序列$x$, 需要求解
\[
y^* = \arg \max_y P(y|x)
\]

由于规范化因子$Z(x)$只依赖于$x$, 与$y$无关, 因此解码问题等价于最大化未归一化得分:
\[
y^*
    = \arg \max_y \sum_{i=1}^n \sum_k \omega_k f_k (y_{i-1}, y_i, x, i)
\]

\subsubsection{维特比算法}


线性链CRF的最优路径可以用维特比算法求解.

定义:
\[
\delta_i (y_i)
    = \max_{y_1, \dots, y_{i-1}}
      \prod_{j=1}^i M_j (y_{j-1}, y_j | x)
\]

写成列向量形式:
\[
\delta_i = \begin{bmatrix}\delta_i(y_1) \\ \delta_i(y_2) \\ \vdots \\ \delta_i(y_m)\end{bmatrix}
\]

初始化为:
\[
\delta_1 = \alpha_1
\]

对$i \ge 2$, 先构造候选矩阵
\[
C_i = M_i^T \operatorname{diag}[\delta_{i-1}]
\]
其中$C_i$第$(r,s)$个元素表示“从前一状态$y_s$转移到当前状态$y_r$时的路径得分”.

则递推公式可以写成矩阵形式:
\[
\delta_i = \operatorname{rowmax}[C_i]
\]

同时记录回溯指针:
\[
\psi_i = \arg\operatorname{rowmax}[C_i]
\]

最后先计算
\[
y_n^* = \arg \max_j \delta_n (y_j)
\]
再沿着$\psi_i$逐步回溯, 得到整条最优路径.

\subsection{模型学习}


CRF的学习本质上是一个凸优化问题, 常用算法包括:

\begin{itemize}
\item 梯度上升法.
\item 拟牛顿法, 如L-BFGS.
\item 其他数值优化方法.
\end{itemize}

\section{HMM和条件随机场}


HMM和CRF都用于序列标注, 但它们的出发点不同.

\begin{table}[H]
\centering
\begingroup
\setlength{\tabcolsep}{2pt}
\small
\begin{tabular}{|p{0.28\textwidth}|p{0.28\textwidth}|p{0.28\textwidth}|}
\hline
比较维度 & HMM & CRF \\ \hline
建模对象 & $P(x, y)$ & $P(y|x)$ \\ \hline
模型类型 & 生成式模型 & 判别式模型 \\ \hline
观测假设 & 当前观测只依赖当前状态 & 可使用更丰富的输入特征 \\ \hline
特征表示 & 通常较固定 & 可灵活设计重叠特征 \\ \hline
归一化方式 & 局部概率归一化 & 全局条件概率归一化 \\ \hline
适用任务 & 生成与识别都可 & 更适合监督式标注任务 \\ \hline
\end{tabular}
\endgroup
\end{table}

两者的主要区别可以概括为:

\begin{itemize}
\item HMM需要对观测序列的生成过程做更强假设.
\item CRF直接对条件概率建模, 能更灵活地融合上下文特征.
\item 在判别任务中, CRF通常比HMM具有更强的表达能力.
\end{itemize}

但与此同时:

\begin{itemize}
\item HMM结构更简单, 推导更直观.
\item CRF训练代价更高, 需要反复计算规范化因子和特征期望.
\end{itemize}

因此, 可以把CRF看成是在序列标注问题上对HMM和一般对数线性模型的一种综合与推广.
